\documentclass[../main.tex]{subfiles}

\begin{document}

(Vázlat)

4.1. Kísérleti setup: A „szándékos hiba”
\begin{itemize}
    \item A feladat: Egy egyszerű API integráció, ahol egy elavult könyvtárverzió miatt a „hallucinált” vagy generált kód futásidejű hibát eredményez.
    \item A hiba jellege: Nem szintaktikai, hanem logikai/környezeti, amit az AI nem lát át közvetlen kontextus nélkül.
\end{itemize}

4.2. Mért változók és mérőszámok
\begin{itemize}
    \item Persistence metrika: Hány egymást követő promptot küld a felhasználó anélkül, hogy manuálisan módosítaná a kódot vagy elolvasná a dokumentációt.
    \item Latency to Manual Shift: Az az idő, ami eltelik az első hibaüzenet és az első manuális debug-lépés (pl. console.log vagy debugger használata) között.
    \item Frusztrációs skála: Likert-skálás kérdőív a „trial-and-error loop” okozta stressz mérésére.
\end{itemize}

\end{document}
